\chapter*{Preface}
\addcontentsline{toc}{chapter}{Preface}
\markboth{Preface}{Preface}

This book is about a deceptively simple question: \emph{how do you fit a
self-exciting process when you never see the individual events, only how many
happened in each time bucket?} Daily hospital admissions, hourly transaction
volumes, per-day video views --- in all of these the natural model is a Hawkes
process (events trigger more events), yet the data are \emph{counts}, not
timestamps, and the Hawkes likelihood needs timestamps. The resolution we build,
piece by piece, is the \emph{Mean Behavior Poisson process} (MBPP): a companion
process whose deterministic intensity equals the expected Hawkes intensity, which
\emph{can} be fit to counts and which shares its parameters with the Hawkes
process we actually want.

\paragraph{Who this is for.} A reader comfortable with calculus, linear algebra,
and a little probability, who wants both the \emph{mathematics} (stated
carefully, with proofs) and the \emph{code} (a working Python package,
\code{hawkes\_calibration}) --- and, above all, the \emph{reasons} behind each
modelling and computational choice. Every chapter explains \emph{why} we do what
we do before showing \emph{how}.

\paragraph{How to read it.} The chapters build strictly in difficulty:

\begin{itemize}[leftmargin=1.4em]
\item \textbf{Part I} assembles the two ingredients --- the Poisson process (and the
one theorem that makes the whole approach possible) and the Hawkes process (and why
it resists counts).
\item \textbf{Part II} constructs the MBPP, solves its governing equation, fits it
to counts, and asks what is actually learnable.
\item \textbf{Part III} adds realism: covariates in the baseline, then --- the hard
case --- covariates in the \emph{excitation}, which turns the governing equation
into a genuine Volterra system.
\item \textbf{Part IV} is computation: the family of exact and numerical solvers,
the operator viewpoint, and learning the solver with neural networks (DeepONet,
FNO, PINN, PINO).
\item \textbf{Part V} is a guided tour of the package and its experiments.
\item The \textbf{appendices} are self-contained primers (Volterra equations,
Laplace/Fourier, ODE/state-space, estimation theory) and an API cheat-sheet.
\end{itemize}

\paragraph{The boxes.} Coloured boxes carry the pedagogy. \emph{Why we care}
motivates; \emph{Intuition} gives the one-sentence mental model; \emph{Why this
choice} defends a design decision against alternatives; \emph{Worked example} does
a concrete computation; \emph{Pitfall} flags a trap; and each chapter ends with a
\emph{recap}. Grey code listings show the actual implementation, and short
exercises invite you to check understanding.

\paragraph{A note on honesty.} We are candid about limitations throughout --- which
parameters are well identified and which are not, where a method is exact and where
it is only an $O(\Delta t)$ approximation, and which models are easy versus hard to
fit. Good modelling is as much about knowing the failure modes as the successes.

\bigskip
\noindent\emph{All code shown is from the companion package; every numerical claim
in the book is reproduced by its test suite or experiments.}
